\cleardoublepage

\section{\textbf{HieraTest}算法}

\subsection{问题描述与设计思路}

本文研究的问题可以表述为：给定待测程序代码及其上下文信息，自动生成一组可编译、可执行且尽可能具有较高覆盖能力的单元测试，并在生成结果出现编译错误或运行错误时，对测试代码进行自动修复。该问题同时涉及代码理解、测试构造、错误定位和结果校验等多个子任务，单一模型往往难以在所有子任务上都保持稳定表现。

基于这一观察，HieraTest 采用大小模型协同的设计思路。大模型主要负责复杂语义理解、关键决策和修复建议生成，小模型主要负责轻量候选生成、局部变体扩展和快速验证。通过将任务拆分为多个阶段，并在不同阶段引入不同能力层级的模型，HieraTest 试图在生成质量、修复效率和推理成本之间取得平衡。

\subsection{方法概述}

HieraTest 的整体流程包括四个阶段：代码分析、测试生成、测试修复和结果筛选。首先，系统对待测类及其依赖进行静态分析，提取方法签名、字段信息、调用关系和可访问上下文。其次，系统根据方法复杂度和可测试性，将生成任务分派给小模型或大模型，生成候选测试代码。然后，系统对候选测试执行编译与运行检查，并收集失败信息；若存在错误，则进入修复阶段，由模型根据错误类型和上下文生成补丁。最后，系统对通过编译和执行验证的测试样例进行保留，并输出最终测试集。

为了更好地体现协同特性，HieraTest 并不将所有任务都交给同一模型完成，而是根据任务粒度和难度进行分层处理：简单、规则性较强的部分交由小模型提升效率，复杂、依赖语义判断的部分交由大模型提升准确率。

\subsection{静态分析模块}

静态分析模块负责从目标项目中提取生成测试所需的基础信息，包括类路径、包名、方法列表、参数类型、返回值类型、异常声明以及可见字段等。该模块还会结合程序依赖图、调用关系和局部控制流信息，为后续提示词构造提供结构化上下文。

在单元测试生成任务中，静态分析的作用不仅在于帮助模型理解目标代码，还在于约束生成空间。例如，当目标方法依赖特定构造函数或辅助方法时，静态分析可以显式提供这些依赖，从而降低模型遗漏导入和初始化逻辑的概率。

\subsection{协同生成模块}

协同生成模块是 HieraTest 的核心。该模块首先根据目标方法的复杂度、参数个数、返回类型和历史失败情况，判断应优先调用大模型还是小模型。对于结构简单、输入输出关系明确的方法，小模型可以快速生成测试候选；对于包含复杂分支、异常路径或对象交互的方法，大模型则更适合生成更完整的测试代码。

在生成过程中，系统还会对候选结果进行轻量筛选，例如检查是否包含必要的断言、是否正确引用被测类、是否满足测试框架约束等。若候选数量不足，系统可以继续调用另一层级模型补充生成，从而形成互补式生成机制。

\subsection{协同修复模块}

当候选测试在编译或执行阶段失败时，HieraTest 会启动协同修复流程。该流程首先根据错误信息识别失败类型，例如导入缺失、类型不匹配、断言错误、对象初始化错误和运行时异常等；然后根据错误类型将修复任务分配给合适的模型。小模型可用于快速尝试局部修改，大模型则用于处理复杂错误场景或生成更稳健的修复建议。

修复后，系统重新执行编译与测试验证，并根据结果决定是否继续迭代。通过引入错误反馈闭环，HieraTest 可以在一定程度上提升测试样例的最终可执行性。

\subsection{结果输出与评估接口}

HieraTest 的输出不仅包括最终通过验证的测试代码，还包括中间生成过程中的错误日志、修复记录和覆盖率数据。为了便于实验评估，系统将生成结果组织为统一目录结构，并输出可供后续统计的文件。这样既方便人工检查样例，也方便统计编译通过率、运行通过率和覆盖率等指标。
